% Build from the repository root (auxiliary files stay under build/):
% mkdir -p build/camera-tree-writeup
% latexmk -pdf -interaction=nonstopmode -halt-on-error \
%   -outdir=build/camera-tree-writeup docs/technical/schur_complement_merging.tex
% cp build/camera-tree-writeup/schur_complement_merging.pdf docs/technical/
\documentclass[11pt,letterpaper]{article}
\usepackage[T1]{fontenc}
\usepackage{lmodern}
\usepackage[margin=0.9in]{geometry}
\usepackage{amsmath,amssymb,bm,mathtools}
\usepackage{booktabs,tabularx}
\usepackage{microtype}
\usepackage{xcolor}
\usepackage{enumitem}
\usepackage{fancyhdr}
\usepackage[colorlinks=true,linkcolor=blue!45!black,urlcolor=blue!45!black]{hyperref}
\hypersetup{pdftitle={Merging Reconstructions with Landmark Schur Complements},pdfauthor={GTSFM technical note}}
\pagestyle{fancy}
\fancyhf{}
\fancyhead[L]{\small GTSFM technical note}
\fancyhead[R]{\small Camera-tree merging}
\fancyfoot[C]{\thepage}
\setlength{\headheight}{14pt}
\setlength{\parindent}{0pt}
\setlength{\parskip}{5pt plus 1pt}
\setlist{nosep,leftmargin=*}
\newcommand{\R}{\mathbb{R}}
\newcommand{\SE}{\mathrm{SE}(3)}
\newcommand{\Sim}{\mathrm{Sim}(3)}
\newcommand{\Local}{\operatorname{Local}}
\newcommand{\rank}{\operatorname{rank}}
\newcommand{\diag}{\operatorname{diag}}
\newcommand{\argmin}{\operatorname*{arg\,min}}
\newcommand{\code}[1]{\texttt{\detokenize{#1}}}
\begin{document}
\thispagestyle{empty}
\begin{center}
{\LARGE\bfseries Merging Reconstructions with\\[4pt] Landmark Schur Complements\par}
\vspace{9pt}
{\large Joint camera poses and intrinsics across local coordinate frames}
\end{center}
\vspace{3pt}
\begin{abstract}
\noindent Independently reconstructed image clusters can be merged by summarizing
each cluster as a joint likelihood over its cameras. We linearize the local
reprojection problem and eliminate only landmark increments, retaining camera
poses, intrinsics, and their correlations. Each resulting Gaussian system is
factored into a Bayes tree and expressed in a common frame through a leaf-specific
similarity transformation. Shared images identify common pose and calibration
variables, and one root optimization reconciles all leaf evidence. This note
derives the Schur reduction, coordinate constraints, similarity-aware residuals,
analytic derivatives, and final camera information, and states the assumptions
under which the frozen approximation is useful.
\end{abstract}

\section{Local reconstructions and shared camera variables}
The input is a hierarchy of reconstructed image clusters. Each leaf supplies
camera poses, intrinsics, image measurements, and 3D landmarks in its own
coordinate frame. Overlapping leaves can contain the same image but assign it
different local indices and different estimates. Image filenames establish
identity across leaves. The merged state contains one global pose and calibration
per image and one similarity transformation per leaf.

\begin{center}
\begin{tabularx}{\linewidth}{@{}lX@{}}
\toprule
Symbol & Meaning \\
\midrule
$\ell$, $i$, $j$ & Leaf, globally identified image, and leaf-local landmark indices. \\
$\bar T_{\ell i}\in\SE$ & Frozen camera-to-leaf pose estimate. \\
$\bar\kappa_{\ell i}\in\R^3$ & Frozen \code{Cal3Bundler} parameters $(f,k_1,k_2)$. \\
$\bar p_{\ell j}\in\R^3$ & Frozen landmark position in the leaf frame. \\
$T_i,\kappa_i$ & Global camera pose and per-image calibration to optimize. \\
$S_\ell\in\Sim$ & Similarity mapping global coordinates into leaf $\ell$. \\
\bottomrule
\end{tabularx}
\end{center}

The observation model uses one focal length and two radial-distortion coefficients.
For $q=T^{-1}p$, let $x=q_x/q_z$, $y=q_y/q_z$, and $r^2=x^2+y^2$. Then
\begin{equation}
 \pi(T,\kappa,p)=
 \begin{bmatrix} f(1+k_1r^2+k_2r^4)x+u_0\\
                 f(1+k_1r^2+k_2r^4)y+v_0\end{bmatrix}.
\end{equation}
The principal point $(u_0,v_0)$ is fixed. Occurrences of a shared image must agree
on it. Different images have independent calibrations, while occurrences of the
same image share one global calibration.

Each leaf's reprojection model is linearized once at its supplied estimates.
The reconstruction-based model uses isotropic pixel noise $\sigma^2 I_2$,
with $\sigma=1$ by default, and excludes tracks observed by fewer than two distinct
retained cameras. The resulting camera likelihood retains the weights of the
linearized model. After landmark elimination, merging uses only that likelihood,
the frozen camera estimates, and image identities; it does not revisit the points.

\newpage
\section{Eliminating landmarks while retaining intrinsics}
For a leaf, collect pose increments $\delta x\in\R^{6n}$, calibration increments
$\delta\kappa\in\R^{3n}$, and landmark increments $\delta p\in\R^{3m}$.
Whitened linearization gives the local least-squares problem
\begin{equation}\label{eq:linear}
 \min_{\delta x,\delta\kappa,\delta p}
 \frac12\left\|A_x\delta x+A_\kappa\delta\kappa+A_p\delta p-b\right\|^2,
 \qquad b=-r(\bar T,\bar\kappa,\bar p).
\end{equation}
The residual offset $b$ is retained even when the supplied reconstruction is not
at a stationary point. Set $\delta c=[\delta x^\top\;\delta\kappa^\top]^\top$
and $A_c=[A_x\;A_\kappa]$. The normal equations have block structure
\begin{equation}
 \begin{bmatrix}H_{cc}&H_{cp}\\H_{pc}&H_{pp}\end{bmatrix}
 \begin{bmatrix}\delta c\\\delta p\end{bmatrix}
 =\begin{bmatrix}g_c\\g_p\end{bmatrix},\qquad
 H=A^\top A,\quad g=A^\top b.
\end{equation}
If $H_{pp}$ is nonsingular, profiling out the points yields
\begin{align}
 H_{\mathrm{red}}&=H_{cc}-H_{cp}H_{pp}^{-1}H_{pc},\label{eq:schur}\\
 g_{\mathrm{red}}&=g_c-H_{cp}H_{pp}^{-1}g_p.\label{eq:linear-term}
\end{align}
The pose--intrinsics block is therefore
\begin{equation}
 (H_{\mathrm{red}})_{x\kappa}
 =H_{x\kappa}-H_{xp}H_{pp}^{-1}H_{p\kappa}.
\end{equation}
Preserving this block allows calibration changes to influence the pose solution
and preserves correlations between images induced by shared landmarks.

\subsection{Square-root reduction}
The implementation uses QR elimination rather than explicitly forming
Equation~\eqref{eq:schur}. Conceptually, an orthogonal transformation gives
\begin{equation}
 Q^\top\begin{bmatrix}A_p&A_c&b\end{bmatrix}
 =\begin{bmatrix}R_p&B_p&d_p\\0&B&d\end{bmatrix}.
\end{equation}
When the landmark columns are eliminable, solving the first block for $\delta p$
leaves the reduced cost
\begin{equation}\label{eq:reduced}
 E_{\mathrm{red}}(\delta c)=\tfrac12\|B\delta c-d\|^2,
 \qquad B^\top B=H_{\mathrm{red}},\quad B^\top d=g_{\mathrm{red}}.
\end{equation}
Any additive constant discarded by the factor representation does not affect
optimization. GTSAM's \code{GaussianFactorGraph.marginal(retained_keys, EliminateQR)}
retains the pose and calibration keys and eliminates only landmark keys.
The returned producer summary is ungauged.

\subsection{The retained camera likelihood}
The reduced state has $9n$ coordinates: six pose and three calibration coordinates
per image. The retained likelihood includes pose--pose, calibration--calibration,
and pose--calibration interactions. Replacing it with independent priors on each
camera would generally lose information about how cameras can change together.
The linear term is equally essential: $g_{\mathrm{red}}$ determines the preferred
increment when the observations have nonzero residuals. An information matrix
alone is insufficient, since $x^2$ and $(x-3)^2$ have the same quadratic coefficient
but different optima. The summary therefore preserves both $B$ and $d$, or an
equivalent collection of Gaussian factors, together with its linearization point.

\newpage
\section{Gauge constraints and leaf Bayes trees}
An observable monocular leaf has seven similarity gauge freedoms: three
rotations, three translations, and one scale. Intrinsics remain in the state,
so for $n$ cameras the expected ungauged rank is $9n-7$.
Additional self-calibration or geometric ambiguities are rejected; the code does
not add implicit calibration priors to make such a system invertible.

The consumer fixes a local coordinate chart before factoring an ungauged leaf.
It adds six exact rows fixing the first pose at its frozen estimate and one exact
row fixing the baseline length to the farthest camera, linearized at those
estimates. A nonzero baseline is required. These constraints act only on poses,
not on calibration variables. GTSAM's \code{NonlinearEqualityPose3} and constrained
\code{RangeFactorPose3} supply the rows.

The augmented graph is factored by multifrontal QR with GTSAM's default ordering.
Its Gaussian Bayes-tree conditionals jointly represent the constrained camera
quadratic. A conditional can contain any mixture of pose and calibration blocks;
block widths are six and three, respectively. Conditionals preserve their linear
terms and noise models, including exact constraints.

A supplied camera summary may already include coordinate constraints. In that
case its retained system must be full rank, and no second gauge is added.
Gauged and ungauged leaves can coexist in the hierarchy. Rank checks normalize
Jacobian columns and tolerate numerical perturbations of gauge null modes.
The constraints must express a choice of coordinates rather than unintended
geometric priors; rank alone cannot distinguish these meanings.

\section{Lifting the leaf quadratic to global variables}
A similarity variable lets each leaf's constrained local chart move within the
global reconstruction. Let $\Phi(S,T)$ denote GTSAM's
\code{Similarity3.transformFrom(Pose3)}. With GTSAM's stored parameters
$S=(R_s,t_s,s)$, its action is
\begin{equation}
 p\longmapsto s(R_sp+t_s),\qquad
 \Phi(S,(R,t))=(R_sR,\;s(R_st+t_s)).
\end{equation}
The translation parameter is inside the scale in this convention. The code uses
the library operation and its derivatives, rather than implementing an alternate
similarity convention.

For each image in leaf $\ell$, define frozen-chart increments
\begin{align}
 \xi_{\ell i}(T_i,S_\ell)
 &=\Local_{\bar T_{\ell i}}\!\left(\Phi(S_\ell,T_i)\right)\in\R^6,\\
 \eta_{\ell i}(\kappa_i)
 &=\Local_{\bar\kappa_{\ell i}}(\kappa_i)
 =\kappa_i-\bar\kappa_{\ell i}\in\R^3.
\end{align}
For conditional $q$, concatenate the appropriate increments in its actual key
order into $z_{\ell q}$. The lifted residual is
\begin{equation}\label{eq:lift}
 e_{\ell q}=A_{\ell q}z_{\ell q}-b_{\ell q}.
\end{equation}
Each shared image has one global calibration but may have different frozen
calibration estimates in different leaves. Similarities affect poses only:
neither focal length nor distortion coefficients change under a change of world
frame or world scale. The residual is passed unwhitened to a \code{CustomFactor}
with the conditional's original noise model, so weighting is applied exactly once.

\newpage
\section{Derivatives, root optimization, and the final reduction}
The lifted factors use analytic derivatives in GTSAM's local retraction charts.
For a pose block, let $A_i^x$ be its six columns in a conditional and define
\begin{equation}
 J_i^{\mathrm{loc}}=D_2\Local_{\bar T_{\ell i}}(\Phi(S_\ell,T_i)),\quad
 J_i^S=D_S\Phi(S_\ell,T_i),\quad J_i^T=D_T\Phi(S_\ell,T_i).
\end{equation}
The chain rule gives
\begin{align}
 \frac{\partial e_{\ell q}}{\partial \delta T_i}
   &=A_i^xJ_i^{\mathrm{loc}}J_i^T,\\
 \frac{\partial e_{\ell q}}{\partial \delta S_\ell}
   &=\sum_{i\text{ in pose blocks}}A_i^xJ_i^{\mathrm{loc}}J_i^S,\\
 \frac{\partial e_{\ell q}}{\partial \delta\kappa_i}
   &=A_i^\kappa.
\end{align}
Calibration contributes no similarity derivative. In particular, a
calibration-only conditional has an identically zero similarity Jacobian.
The implementation computes column offsets from the variable types; it never
assumes that every conditional block has six columns. These formulas apply at
nonidentity poses and nonunit similarity scales, not just at initialization.

\subsection{Initialization and one global solve}
Camera identity is resolved by filename, and sorted filenames determine the
output indices. Pose-based similarity alignment initializes a connected graph of
overlapping leaves. Each usable overlap link needs at least two shared cameras
with distinct centers. Global calibrations are initialized from the first
occurrence of each image in leaf traversal order.

The input scene tree assembles the leaf evidence bottom-up, without performing
an optimization at internal nodes. The root problem is
\begin{equation}\label{eq:root}
 \min_{\{T_i,\kappa_i,S_\ell\}}\ \sum_{\ell,q}
   \tfrac12\|e_{\ell q}(T,\kappa,S_\ell)\|_{\Sigma_{\ell q}}^2,
 \qquad S_0=I,
\end{equation}
Here $\|e\|_\Sigma^2=e^\top\Sigma^{-1}e$ for stochastic rows; zero-noise rows
are exact constraints. Fixing the first leaf's similarity selects the global frame and
scale. Levenberg--Marquardt with multifrontal QR solves this problem once at the
root. The local pose and calibration estimates used inside the factors remain
frozen throughout optimization. The public leaf-to-global transform is
$S_\ell^{-1}$.

\subsection{Final camera information retains calibration}
At the converged solution, the assembled graph is linearized again with respect
to its global camera and similarity variables. This relinearizes the lifted
factors, not the original image measurements. The final marginalization removes
only the similarity increments:
\begin{equation}
 E_{\mathrm{final}}(\delta c)
 =\min_{\delta s} E_{\mathrm{root,lin}}(\delta c,\delta s),\qquad
 \delta c=(\delta x,\delta\kappa).
\end{equation}
The resulting Gaussian Bayes tree contains global \code{X(i)} pose keys and
\code{K(i)} calibration keys. It contains no landmarks or similarities and
includes the chosen coordinate gauge. The returned scene has optimized poses
and calibrations, image metadata, and zero tracks. Its quadratic describes
local information in the frozen approximation, not the uncertainty of a newly
optimized full bundle-adjustment problem.

\newpage
\section{The complete merging procedure}
The algorithm separates measurement reduction from camera merging. A leaf
summary contains the reduced Gaussian factors, frozen pose and calibration
estimates, image identities, local key associations, and a declaration of whether
its coordinate gauge is fixed. The global optimizer needs no landmark positions
or image observations once these summaries have been constructed.

\begin{enumerate}[itemsep=2pt]
\item \textbf{Linearize each leaf.} Assemble the weighted local observation model
at the supplied poses, intrinsics, and points, retaining the residual offsets.
\item \textbf{Eliminate landmarks.} Apply QR to the point columns and retain all
pose and calibration columns. Preserve the joint quadratic and linear term.
\item \textbf{Choose local coordinates and factor.} For each ungauged leaf, fix
one pose and one baseline length. Factor the retained graph into Gaussian
Bayes-tree conditionals.
\item \textbf{Match images and initialize frames.} Assign one global pose and
calibration to each filename. Align overlapping leaves to initialize their
similarities, and initialize calibration from the first occurrence of each image.
\item \textbf{Lift and assemble.} Evaluate each conditional through its leaf's
similarity and frozen camera charts. Collect the leaf factors bottom-up, counting
each leaf's contribution once.
\item \textbf{Optimize globally.} Fix the first similarity to identity and jointly
refine all global poses, intrinsics, and remaining similarities at the root.
\item \textbf{Return cameras and local information.} Export the optimized cameras
and leaf-to-global similarities. Linearize the merged graph and remove the
similarity increments to obtain the final joint camera Bayes tree.
\end{enumerate}

\subsection{The scene tree and the Bayes trees have different roles}
The scene tree organizes which reconstructions contribute evidence and how their
summaries are collected. Internal nodes concatenate child factor graphs; they do
not optimize or discard additional camera variables. Each leaf's Bayes tree is
an algebraic factorization of that leaf's reduced likelihood, and the final Bayes
tree factors the merged camera information. Their clique structure follows
variable connectivity and elimination ordering, not the scene hierarchy.

\subsection{Assumptions and approximation limits}
The usable overlap graph must be connected, with at least two shared cameras at
distinct centers on each link used for initialization. Each leaf must have the
declared geometric rank; additional self-calibration ambiguity is not repaired
by implicit calibration priors. Image identities and fixed principal points must
agree across shared images. Evidence must be assigned consistently across leaves:
sharing camera variables is intentional, but counting the same observation twice
would duplicate its information.

Landmark elimination is exact for the linearized problem, up to additive constants.
During merging, pose-chart and similarity derivatives change but the leaf
quadratics remain fixed. Large departures from their linearization points can
reduce accuracy. Recovering landmarks requires a subsequent step using the
original observations. Uncertainty interpretation must respect the chosen gauge
and the frozen-linearization approximation.
\end{document}
